% ===================================================================
%  ਸਾਰਣੀ ਨੰ. 2 — ਮਨੁੱਖੀ ਸਰੀਰ। — ਛੁੱਟੀ। — ਖੇਡਾਂ।
%  Traduction 2026 d'apres texto/io, controlee sur texto/fr, texto/en
%  et texto/hi. Consigne : voir texto/pa/10-sarni-01.tex.
% ===================================================================

%%K t02-apar-1 apar
\VUsaut{4.0mm}
\VUtitre{15.0pt}{60}{ਸਾਰਣੀ ਨੰ. 2}
\VUsaut{2.0mm}
\VUtitre{11.4pt}{0}{ਮਨੁੱਖੀ ਸਰੀਰ। --- ਛੁੱਟੀ।}
\VUtitre{11.4pt}{0}{ਖੇਡਾਂ।}

%%K t02-tit-0 sub
\VUtitre{13.2pt}{0}{ਮਨੁੱਖੀ ਸਰੀਰ।}
\VUsaut{2.4mm}
\VUcentre{10.2pt}{0}{\textit{(ਕੁਦਰਤੀ ਵਿਗਿਆਨ ਦਾ ਇਕ ਸੌਖਾ ਪਾਠ।)}}
\VUsaut{3.0mm}

%%K t02-01-1 p
1. --- ਮਨੁੱਖੀ ਸਰੀਰ ਦੇ ਮੁੱਖ ਹਿੱਸੇ ਇਹ ਹਨ: \VUgras{ਸਿਰ}, \VUgras{ਧੜ} ਤੇ
\VUgras{ਅੰਗ}।

%%K t02-tit-1 sub
\VUpk{10.2pt}{40}{ਸਿਰ।}
\VUsaut{3.0mm}

%%K t02-02-1 p
2. --- ਸਿਰ ਦੇ ਦੋ ਹਿੱਸੇ ਹਨ: \VUgras{ਖੋਪੜੀ}, ਜਿਸ ਵਿਚ \VUgras{ਦਿਮਾਗ਼} ਰਹਿੰਦਾ ਹੈ
ਤੇ ਜਿਸ ਨੂੰ \VUgras{ਵਾਲ} ਢਕੀ ਰੱਖਦੇ ਹਨ, ਤੇ \VUgras{ਮੂੰਹ}।

%%K t02-03-1 p
3. --- ਸਾਡੇ ਚਿੱਤਰ ਵਿਚ ਨਾ ਖੋਪੜੀ ਵਿਖਾਈ ਦਿੰਦੀ ਹੈ ਨਾ ਦਿਮਾਗ਼; ਪਰ ਮੂੰਹ ਦੇ ਮੁੱਖ ਹਿੱਸੇ
ਸਾਨੂੰ ਬਹੁਤ ਸਾਫ਼ ਵਿਖਾਈ ਦਿੰਦੇ ਹਨ।

%%K t02-04-1 p
4. --- ਸਭ ਤੋਂ ਪਹਿਲਾਂ ਇਹ \VUgras{ਮੱਥਾ} ਹੈ, ਜੋ ਜ਼ੋਰਦਾਰ ਬੁੱਧੀ ਤੇ ਸਦਾ ਕੰਮ ਕਰਦੇ ਮਨ
ਦਾ ਪਤਾ ਦਿੰਦਾ ਹੈ। \VUgras{ਪੁੜਪੁੜੀਆਂ} ਬਹੁਤ ਖੁੱਲ੍ਹੀਆਂ ਹਨ। \VUgras{ਅੱਖ} ਚੰਚਲ ਤੇ
ਚੌੜੀ ਖੁੱਲ੍ਹੀ ਹੈ। ਇਕ ਸੰਘਣੀ \VUgras{ਭਰਵੱਟੀ} ਤੇ ਲੰਮੀਆਂ \VUgras{ਪਲਕਾਂ} ਉਸ ਦੀ
ਰਾਖੀ ਕਰਦੀਆਂ ਹਨ।

%%K t02-05-1 p
5. --- ਇੱਥੇ \VUgras{ਨੱਕ} ਤੇ \VUgras{ਨਾਸਾਂ} ਵੀ ਹਨ। ਨੱਕ ਨਾਲ ਅਸੀਂ ਸੁੰਘਦੇ ਤੇ ਸਾਹ
ਲੈਂਦੇ ਹਾਂ। ਨੱਕ ਦੇ ਦੋਹਾਂ ਪਾਸੇ \VUgras{ਗੱਲ੍ਹਾਂ} ਹਨ, ਤੇ ਉਨ੍ਹਾਂ ਤੋਂ ਪਰੇ
\VUgras{ਕੰਨ}। ਮੂੰਹ ਦਾ ਹੇਠਲਾ ਹਿੱਸਾ \VUgras{ਮੁਖ} ਤੇ \VUgras{ਠੋਡੀ} ਨਾਲ ਬਣਦਾ ਹੈ।

%%K t02-06-1 p
6. --- ਉਹ ਸਾਨੂੰ ਚੰਗੀ ਤਰ੍ਹਾਂ ਵਿਖਾਈ ਨਹੀਂ ਦਿੰਦੇ, ਕਿਉਂਕਿ \VUgras{ਮੁੱਛਾਂ} ਤੇ
\VUgras{ਦਾੜ੍ਹੀ} ਉਨ੍ਹਾਂ ਨੂੰ ਲੁਕਾਈ ਰੱਖਦੀਆਂ ਹਨ। ਪਰ ਅਸੀਂ ਜਾਣਦੇ ਹਾਂ ਕਿ ਮੁਖ ਵਿਚ ਦੋ
\VUgras{ਬੁੱਲ੍ਹ}, \VUgras{ਉੱਪਰਲਾ ਜਬਾੜਾ} ਤੇ \VUgras{ਹੇਠਲਾ ਜਬਾੜਾ} ਆਪਣੇ
\VUgras{ਦੰਦਾਂ} ਸਮੇਤ, ਤੇ \VUgras{ਜੀਭ} ਹੁੰਦੀ ਹੈ। ਮੁਖ ਨਾਲ ਅਸੀਂ ਬੋਲਦੇ ਤੇ ਖਾਂਦੇ
ਹਾਂ। ਜੀਭ ਨਾਲ ਅਸੀਂ ਆਪਣੇ ਖਾਣੇ ਦਾ ਸੁਆਦ ਲੈਂਦੇ ਹਾਂ।

%%K t02-07-1 p
7. --- ਅੱਖ, ਕੰਨ, ਨੱਕ ਤੇ ਮੁਖ, ਉਂਗਲਾਂ ਸਮੇਤ, ਪੰਜਾਂ ਗਿਆਨ-ਇੰਦਰੀਆਂ ਦੇ ਅੰਗ ਹਨ:
\VUgras{ਦ੍ਰਿਸ਼ਟੀ}, \VUgras{ਸ੍ਰਵਣ}, \VUgras{ਗੰਧ}, \VUgras{ਸੁਆਦ} ਤੇ
\VUgras{ਸਪਰਸ਼}।

%%K t02-08-1 p
8. --- ਇਸ ਤਰ੍ਹਾਂ ਸਿਰ ਮਨੁੱਖੀ ਸਰੀਰ ਦੇ ਸਭ ਤੋਂ ਦਿਲਚਸਪ ਹਿੱਸਿਆਂ ਵਿਚੋਂ ਇਕ ਹੈ।

%%K t02-tit-2 sub
\VUsaut{4.4mm}
\VUpk{10.2pt}{40}{ਧੜ।}
\VUsaut{3.0mm}

%%K t02-09-1 p
9. --- \VUgras{ਗਰਦਨ} ਸਿਰ ਨੂੰ ਧੜ ਨਾਲ ਜੋੜਦੀ ਹੈ। ਉਸ ਵਿਚ \VUgras{ਗਲਾ} ਤੇ
\VUgras{ਧੌਣ} ਆਉਂਦੇ ਹਨ।

%%K t02-10-1 p
10. --- ਧੜ ਵਿਚ ਵੀ ਬਹੁਤ ਸਾਰੇ ਲੋੜੀਂਦੇ ਅੰਗ ਰਹਿੰਦੇ ਹਨ। \VUgras{ਛਾਤੀ} ਵਿਚ
\VUgras{ਫੇਫੜੇ} ਹਨ, ਜਿਨ੍ਹਾਂ ਨਾਲ ਅਸੀਂ ਸਾਹ ਲੈਂਦੇ ਹਾਂ, ਤੇ \VUgras{ਦਿਲ}, ਜੋ ਜੀਵਨ
ਦਾ ਕੇਂਦਰ ਹੈ। ਜਦੋਂ ਦਿਲ ਧੜਕਣੋਂ ਹਟ ਜਾਂਦਾ ਹੈ, ਜੀਵ ਮਰ ਜਾਂਦਾ ਹੈ।

%%K t02-11-1 p
11. --- \VUgras{ਪਸਲੀਆਂ} ਛਾਤੀ ਨੂੰ ਸੰਭਾਲਦੀਆਂ ਹਨ।

%%K t02-12-1 p
12. --- ਧੜ ਦੇ ਦੂਜੇ ਹਿੱਸੇ ਹਨ \VUgras{ਪਾਸੇ}, \VUgras{ਪਿੱਠ}, \VUgras{ਮੋਢੇ} ਤੇ
\VUgras{ਢਿੱਡ}। ਅਸੀਂ ਸੌਣ ਲਈ ਪਾਸੇ ਜਾਂ ਪਿੱਠ ਭਾਰ ਲੇਟਦੇ ਹਾਂ, ਤੇ ਭਾਰ ਮੋਢੇ ਉੱਤੇ
ਢੋਂਦੇ ਹਾਂ।

%%K t02-13-1 p
13. --- ਢਿੱਡ ਵਿਚ \VUgras{ਮਿਹਦਾ} ਤੇ \VUgras{ਆਂਦਰਾਂ} ਹਨ। ਉਹ ਸਾਨੂੰ ਖਾਣਾ ਪਚਾਉਣ
ਵਿਚ ਕੰਮ ਆਉਂਦੀਆਂ ਹਨ।

%%K t02-tit-3 sub
\VUsaut{4.4mm}
\VUpk{10.2pt}{40}{ਅੰਗ।}
\VUsaut{3.0mm}

%%K t02-14-1 p
14. --- ਸਾਡੇ ਚਾਰ \VUgras{ਅੰਗ} ਹਨ: ਦੋ \VUgras{ਬਾਹਾਂ} ਤੇ ਦੋ \VUgras{ਲੱਤਾਂ}।
ਬਾਹਾਂ ਨਾਲ ਅਸੀਂ ਚੀਜ਼ਾਂ ਲੈਂਦੇ, ਫੜਦੇ, ਚੁੱਕਦੇ ਤੇ ਇਕੱਠੀਆਂ ਕਰਦੇ ਹਾਂ; ਲੱਤਾਂ ਨਾਲ
ਅਸੀਂ ਖੜੇ ਹੁੰਦੇ, ਤੁਰਦੇ ਤੇ ਦੌੜਦੇ ਹਾਂ।

%%K t02-15-1 p
15. --- \VUgras{ਮੋਢਾ} ਬਾਂਹ ਨੂੰ ਸਰੀਰ ਨਾਲ ਜੋੜਦਾ ਹੈ। ਇਕ ਬਹੁਤ ਮਜ਼ਬੂਤ
\VUgras{ਪੱਠਾ}, ਦੋ-ਸਿਰਾ, ਬਾਂਹ ਨੂੰ ਹਿਲਾਉਂਦਾ ਹੈ, ਜਿਸ ਨੂੰ \VUgras{ਅਰਕ}
\VUgras{ਬਾਂਹ ਦੇ ਹੇਠਲੇ ਹਿੱਸੇ} ਨਾਲ ਜੋੜਦੀ ਹੈ। \VUgras{ਗੁੱਟ} \VUgras{ਹੱਥ} ਦਾ ਜੋੜ
ਬਣਾਉਂਦਾ ਹੈ, ਜੋ \VUgras{ਉਂਗਲਾਂ} ਤੇ \VUgras{ਨਹੁੰਆਂ} ਉੱਤੇ ਮੁੱਕਦਾ ਹੈ। ਹੱਥ ਉੱਤੇ
ਸਾਨੂੰ ਇਕ \VUgras{ਨਾੜ} (ਤੇ ਇਕ ਧਮਣੀ) ਵਿਖਾਈ ਦਿੰਦੀ ਹੈ, ਜਿਸ ਵਿਚ \VUgras{ਲਹੂ} ਵਗਦਾ
ਹੈ।

%%K t02-16-1 p
16. --- ਲੱਤ ਦੇ ਹਿੱਸੇ ਇਹ ਹਨ: \VUgras{ਪੱਟ}, \VUgras{ਗੋਡਾ} ਤੇ
\VUgras{ਗੋਡੇ ਦਾ ਪਿਛਲਾ ਪਾਸਾ}, \VUgras{ਪਿੰਜਣੀ}, ਜਿਸ ਦਾ \VUgras{ਮਾਸ}
\VUgras{ਚਮੜੀ} ਨਾਲ ਢਕਿਆ ਹੈ, \VUgras{ਗਿੱਟਾ} ਤੇ \VUgras{ਪੈਰ}। ਲੱਤ ਦੀਆਂ
\VUgras{ਹੱਡੀਆਂ} ਬਹੁਤ ਮੋਟੀਆਂ ਤੇ ਬਹੁਤ ਪੱਕੀਆਂ ਹਨ। ਪੈਰ ਦੇ ਸਿਰੇ ਉੱਤੇ
\VUgras{ਪੈਰ ਦੀਆਂ ਉਂਗਲਾਂ} ਹਨ। ਦੂਜੇ ਹਿੱਸੇ \VUgras{ਤਲੀ} ਤੇ \VUgras{ਅੱਡੀ} ਹਨ।

%%K t02-17-1 p
17. --- ਮਨੁੱਖੀ ਸਰੀਰ ਦਾ ਤਕਰੀਬਨ ਪੂਰਾ ਵੇਰਵਾ ਇਹੋ ਹੈ।

%%K t02-17-2 p
\textit{ਟਿੱਪਣੀ।} --- \textit{« ਮੇਰਾ ... ਦੁਖਦਾ ਹੈ (ਮੇਰਾ ਸਿਰ, ਆਦਿ) » ਵਾਕ ਦੀ ਮਦਦ
ਨਾਲ ਅਧਿਆਪਕ ਇਹ ਸਾਰੀ ਸ਼ਬਦਾਵਲੀ ਚੰਗੀ ਜਾਂ ਮਾੜੀ} \VUgras{ਸਰੀਰਕ ਸਿਹਤ}
\textit{ਦੇ ਪੱਖੋਂ ਫਿਰ ਦੁਹਰਾ ਸਕਣਗੇ।}

%%K t02-tit-4 sub
\VUsaut{5.0mm}
\VUpk{10.2pt}{40}{ਕਸਰਤ।}
\VUsaut{2.4mm}
\VUcentre{10.2pt}{0}{\textit{(ਇਕ ਵਿਦਿਆਰਥੀ ਦੀ ਜ਼ਬਾਨੀ।)}}
\VUsaut{3.0mm}

%%K t02-18-1 p
18. --- ਲਚਕੀਲੇ, ਫੁਰਤੀਲੇ ਤੇ ਤੰਦਰੁਸਤ ਰਹਿਣ ਲਈ ਸਾਨੂੰ ਆਪਣੀਆਂ ਲੱਤਾਂ ਤੇ ਬਾਹਾਂ ਤੋਂ
ਕੰਮ ਲੈਣਾ ਚਾਹੀਦਾ ਹੈ। ਇਸੇ ਲਈ ਅਸੀਂ ਖੇਡਾਂ ਖੇਡਦੇ ਹਾਂ ਤੇ \VUgras{ਕਸਰਤ} ਕਰਦੇ ਹਾਂ।

%%K t02-19-1 p
19. --- ਸਾਰਾ ਹਫ਼ਤਾ ਇਹ ਦੋਵੇਂ ਅਭਿਆਸ ਸਕੂਲ ਦੇ ਵਿਹੜੇ ਵਿਚ ਹੁੰਦੇ ਹਨ (ਵੇਖੋ ਸਾਰਣੀ ਨੰ.
1)। ਪਰ ਵੀਰਵਾਰ ਤੇ ਐਤਵਾਰ ਨੂੰ ਅਸੀਂ ਇਕ ਵੱਡੇ ਮੈਦਾਨ ਵਿਚ ਜਾਂਦੇ ਹਾਂ, ਜਿੱਥੇ ਦੌੜਨ ਤੇ
ਮੌਜ ਕਰਨ ਦੀ ਥਾਂ ਮਿਲਦੀ ਹੈ।

%%K t02-20-1 p
20. --- ਸਾਡੇ ਅਧਿਆਪਕ ਤੇ ਸਾਡੇ ਮਾਪੇ ਕਦੇ ਕਦੇ ਨਾਲ ਆਉਂਦੇ ਹਨ; ਪਰ ਅਭਿਆਸ
\VUgras{ਕਸਰਤ ਦੇ ਅਧਿਆਪਕ} \textsuperscript{(12)} ਹੀ ਕਰਾਉਂਦੇ ਹਨ।

%%K t02-21-1 p
21. --- ਮੈਦਾਨ ਵਿਚ, ਵੜਨ ਦੇ ਖੱਬੇ ਪਾਸੇ, \VUgras{ਕਸਰਤ ਦੇ ਸਾਧਨ}
\textsuperscript{(1)} ਖੜੇ ਹਨ: ਅਸੀਂ \VUgras{ਪੌੜੀ} \textsuperscript{(2)} ਉੱਤੇ
ਚੜ੍ਹਦੇ ਹਾਂ, \VUgras{ਛੱਲਿਆਂ} \textsuperscript{(3)} ਤੇ \VUgras{ਝੂਲੇ ਦੇ ਡੰਡੇ}
\textsuperscript{(4)} ਉੱਤੇ ਪੁੱਠੇ ਝੂਟੇ ਲੈਂਦੇ ਹਾਂ, \VUgras{ਰੱਸੇ ਦੀ ਪੌੜੀ}
\textsuperscript{(5)} ਜਾਂ \VUgras{ਗੰਢਾਂ ਵਾਲੇ ਰੱਸੇ} \textsuperscript{(6)} ਦੇ ਸਿਰੇ
ਤਕ ਹੱਥੋ-ਹੱਥੀ ਚੜ੍ਹ ਜਾਂਦੇ ਹਾਂ।

%%K t02-22-1 p
22. --- ਇਸ ਦੌਰਾਨ ਸਾਡੇ ਸਾਥੀ \VUgras{ਲੱਕੜ ਦੇ ਘੋੜੇ} \textsuperscript{(7)} ਜਾਂ
\VUgras{ਸਮਾਨੰਤਰ ਡੰਡਿਆਂ} \textsuperscript{(11)} ਦੇ ਉੱਤੋਂ ਟੱਪਦੇ ਹਨ। ਕੁਝ
\VUgras{ਮੁੱਕੇਬਾਜ਼ੀ} \textsuperscript{(8)} ਕਰਦੇ ਹਨ ਜਾਂ ਨਕਾਬ ਤੇ \VUgras{ਫ਼ੌਇਲ}
ਲੈ ਕੇ \VUgras{ਤਲਵਾਰਬਾਜ਼ੀ} \textsuperscript{(9)} ਕਰਦੇ ਹਨ। ਕੁਝ ਅੰਤ ਵਿਚ
\VUgras{ਡੰਬਲ} \textsuperscript{(10)} ਚੁੱਕਦੇ ਹਨ।

%%K t02-tit-5 sub
\VUsaut{4.6mm}
\VUpk{10.2pt}{40}{ਖੇਡਾਂ।}
\VUsaut{3.0mm}

%%K t02-23-1 p
23. --- ਕਸਰਤ ਕਰਨ ਵਾਲਿਆਂ ਦੇ ਚੁਫੇਰੇ ਮੁੰਡੇ ਤੇ ਕੁੜੀਆਂ ਕਈ ਤਰ੍ਹਾਂ ਦੀਆਂ
\VUgras{ਖੇਡਾਂ} ਖੇਡ ਰਹੇ ਹਨ। ਕੁੜੀਆਂ ਆਪਣੇ \VUgras{ਘੇਰੇ} \textsuperscript{(14)} ਨਾਲ
ਦਿਲ ਪਰਚਾਉਂਦੀਆਂ ਹਨ; ਉਹ \VUgras{ਰੱਸੀ ਟੱਪਦੀਆਂ} \textsuperscript{(13)} ਹਨ, ਜਾਂ ਆਪਣੇ
ਭਰਾਵਾਂ ਤੇ ਚਚੇਰੇ ਭਰਾਵਾਂ ਨੂੰ \VUgras{ਕਰੋਕੇ} \textsuperscript{(15)} ਦੀ ਖੇਡ ਵਿਚ
ਸੱਦਦੀਆਂ ਹਨ। ਜਦੋਂ ਵਿਰੋਧੀ ਸੋਚਦਾ ਹੈ ਕਿ ਉਹ ਸੌਖਾ ਹੀ ਘੇਰੇ ਵਿਚੋਂ ਨਿਕਲ ਜਾਵੇਗਾ, ਤਦੇ
ਉਹ ਆਪਣੇ \VUgras{ਲੱਕੜ ਦੇ ਮੁੰਗਲੇ} ਨਾਲ ਉਸ ਦੀ \VUgras{ਗੇਂਦ} ਪਰੇ ਮਾਰ ਦਿੰਦੀਆਂ ਹਨ,
ਤੇ ਇਸ ਨਾਲ ਉਨ੍ਹਾਂ ਨੂੰ ਬਹੁਤ ਮਜ਼ਾ ਆਉਂਦਾ ਹੈ!

%%K t02-24-1 p
24. --- ਮੁੰਡੇ ਵੱਧ ਜ਼ੋਰ ਵਾਲੀਆਂ ਖੇਡਾਂ ਪਸੰਦ ਕਰਦੇ ਹਨ। ਉਹ ਖ਼ੁਸ਼ੀ ਨਾਲ
\VUgras{ਨੌਂ ਗੋਲੀਆਂ} \textsuperscript{(16)} ਖੇਡਦੇ ਹਨ, ਜਿਨ੍ਹਾਂ ਨੂੰ ਉਹ
\VUgras{ਗੇਂਦ} \textsuperscript{(17)} ਨਾਲ ਸਾਫ਼ ਡੇਗ ਦਿੰਦੇ ਹਨ। ਉਹ \VUgras{ਲਾਟੂ}
\textsuperscript{(18)} ਤੇ \VUgras{ਬਿਲਬੋਕੇ} \textsuperscript{(19)} ਨੂੰ, ਇੱਥੋਂ ਤਕ
ਕਿ \VUgras{ਡੱਡੂ ਦੀ ਖੇਡ} \textsuperscript{(20)} ਨੂੰ ਵੀ, ਨਿੱਕਾ ਨਹੀਂ ਸਮਝਦੇ। ਪਰ
ਉਨ੍ਹਾਂ ਦੀਆਂ ਪਿਆਰੀਆਂ ਖੇਡਾਂ \VUgras{ਬੰਟੇ} \textsuperscript{(21)} ਤੇ
\VUgras{ਕੰਧ-ਗੇਂਦ} \textsuperscript{(22)} ਹਨ, ਕਿਉਂਕਿ \VUgras{ਸ਼ੀਸ਼ੇ ਦੇ ਘਰ}
\textsuperscript{(26)} ਦੀ ਕੰਧ ਹੌਲੀ ਗੇਂਦ ਉਛਾਲਣ ਲਈ ਠੀਕ ਬਹਿੰਦੀ ਹੈ।

%%K t02-25-1 p
25. --- ਨਿੱਕਾ ਜਾਨ ਆਪਣੇ \VUgras{ਬਾਂਸਾਂ} \textsuperscript{(24)} ਉੱਤੇ ਚੜ੍ਹਿਆ ਬਹੁਤ
ਸਿਆਣਾ ਹੈ ਤੇ ਆਪਣਾ ਸੰਤੁਲਨ ਖ਼ੂਬ ਸੰਭਾਲਦਾ ਹੈ। ਉਸ ਕੋਲ ਕੁਝ ਸਾਥੀ ਉਸੇ ਸ਼ੀਸ਼ੇ ਦੇ ਘਰ ਵਿਚ
ਤੇ \VUgras{ਵਾੜ} ਦੇ ਪਿੱਛੇ \VUgras{ਲੁਕਣ-ਮੀਟੀ} \textsuperscript{(25)} ਖੇਡ ਰਹੇ ਹਨ।
ਕੁਝ ਨੂੰ \VUgras{ਤਲੀ ਦੀ ਖੇਡ} \textsuperscript{(28)} ਵੱਧ ਚੰਗੀ ਲਗਦੀ ਹੈ, ਜਾਂ
\VUgras{ਚਿੜੀ} \textsuperscript{(29)}, ਜੋ ਇਕ \VUgras{ਬੱਲੇ} ਤੋਂ ਦੂਜੇ ਤਕ ਉੱਡਦੀ
ਰਹਿੰਦੀ ਹੈ।

%%K t02-noto-1 noto
\VUnotes{143.2mm}{%
(*) ਬੱਚਿਆਂ ਦੀਆਂ ਖੇਡਾਂ ਦੇ ਨਾਂ ਅਕਸਰ ਬਿਲਕੁਲ ਕੌਮੀ ਹੁੰਦੇ ਹਨ। ਅਸੀਂ ਉਨ੍ਹਾਂ ਨੂੰ ਇਦੋ
ਵਿਚ ਜਿਵੇਂ ਬਣ ਸਕਿਆ ਉਵੇਂ ਨਾਂ ਦਿੱਤੇ ਹਨ — ਕਦੇ ਉਸ ਕੰਮ ਤੋਂ ਜੋ ਉਨ੍ਹਾਂ ਨੂੰ ਪਛਾਣ
ਦਿੰਦਾ ਹੈ, ਕਦੇ ਕਿਸੇ ਨਾ ਕਿਸੇ ਦੇਸ ਦਾ ਨਾਂ ਅਪਣਾ ਕੇ, ਜਦੋਂ ਉਹ ਬਹੁਤ ਭੁਲੇਖਾ ਪਾਊ ਨਾ
ਹੋਵੇ। ਮਿਸਾਲ ਲਈ: \VUgras{ਪੀਪੇ ਦੀ ਖੇਡ} (ਜਰਮਨ ਤੇ ਫ਼ਰਾਂਸੀਸੀ) ਉਹੋ
\VUgras{ਡੱਡੂ ਦੀ ਖੇਡ} \textsuperscript{(20)} (ਅੰਗਰੇਜ਼ੀ) ਹੈ, ਜਿਸ ਵਿਚ ਖਿਡਾਰੀ ਆਪਣੀਆਂ
ਗੋਲ ਟਿੱਕੀਆਂ ਮੋਰੀਆਂ ਵਾਲੇ ਸੰਦੂਕ (ਪੀਪੇ) ਉੱਤੇ ਮੂੰਹ ਖੋਲ੍ਹੀ ਖੜੇ ਧਾਤ ਦੇ ਡੱਡੂ ਦੇ ਮੂੰਹ
ਵਿਚ ਸੁੱਟਣ ਦਾ ਯਤਨ ਕਰਦੇ ਹਨ। \VUgras{ਮੋਢਾ-ਟੱਪ} \textsuperscript{(30)}
\VUgras{ਬੱਕਰੀ-ਟੱਪ} ਤੋਂ ਸਿਰਫ਼ ਇੰਨੀ ਕੁ ਵੱਖਰੀ ਹੈ: ਸਿਰ ਦੇ ਉੱਤੋਂ ਟੱਪਣ ਲਈ ਖਿਡਾਰੀ
ਆਪਣੇ ਸਾਥੀ ਦੇ ਮੋਢਿਆਂ ਉੱਤੇ ਹੱਥ ਰੱਖਦੇ ਹਨ, ਜਦ ਕਿ « \VUgras{ਬੱਕਰੀ-ਟੱਪ} » ਵਿਚ ਉਹ
ਸਿਰਫ਼ ਉਸ ਦੀ ਪਿੱਠ ਉੱਤੇ ਹੱਥ ਰੱਖਦੇ ਹਨ। --- ਜਾਂ ਫਿਰ, ਕੁਝ ਖਿਡਾਰੀ ਝੁਕ ਜਾਂਦੇ ਹਨ ਤੇ
ਦੂਜੇ ਉਨ੍ਹਾਂ ਦੀ ਪਿੱਠ ਉੱਤੇ ਹੱਥ ਰੱਖ ਕੇ ਉਨ੍ਹਾਂ ਦੇ ਉੱਤੋਂ ਟੱਪਦੇ ਹਨ; ਪਹਿਲਿਆਂ ਨੂੰ ਬਿਨਾਂ
ਡੋਲੇ ਧੱਕਾ ਜਰਨਾ ਪੈਂਦਾ ਹੈ, ਦੂਜਿਆਂ ਨੂੰ ਬਿਨਾਂ ਸੰਤੁਲਨ ਗੁਆਏ ਟੱਪਣਾ (ਸਾਰਣੀ ਹੀ ਵੇਖੋ)।%
}

%%K t02-26-1 p
26. --- ਮੈਦਾਨ ਦੇ ਵਿਚਕਾਰ ਦੋ ਰੁੱਖ ਹਨ। ਉਨ੍ਹਾਂ ਵਿਚੋਂ ਇਕ ਦੇ ਹੇਠਾਂ ਸੱਤ-ਅੱਠ ਮੁੰਡਿਆਂ
ਨੇ \VUgras{ਮੋਢਾ-ਟੱਪ} \textsuperscript{(30)} (*) ਦੀ ਖੇਡ ਜਮਾਈ ਹੋਈ ਹੈ। ਇਹ ਖੇਡ ਹਰ
ਦੇਸ ਵਿਚ ਨਹੀਂ ਜਾਣੀ ਜਾਂਦੀ; ਪਰ \VUgras{ਬੱਕਰੀ-ਟੱਪ} ਕੀ ਹੈ, ਇਹ ਸਾਰੇ ਜਾਣਦੇ ਹਨ,
ਜਿਸ ਨੂੰ ਬੱਚੇ ਇਸ ਵੇਲੇ ਨਹੀਂ ਖੇਡ ਰਹੇ।

%%K t02-tit-6 sub
\VUsaut{5.0mm}
\VUpk{10.2pt}{40}{ਖੁੱਲ੍ਹੀ ਹਵਾ ਦੀਆਂ ਖੇਡਾਂ।}
\VUsaut{3.4mm}

%%K t02-27-1 p
27. --- \VUgras{ਅੱਖ-ਮਿਚੌਲੀ} \textsuperscript{(31)} ਵਿਚ ਵੀ ਬਹੁਤ ਮਜ਼ਾ ਆਉਂਦਾ ਹੈ;
ਕੁੜੀਆਂ ਤੇ ਮੁੰਡੇ ਵਾਰੀ ਵਾਰੀ ਅੱਖਾਂ ਉੱਤੇ \VUgras{ਪੱਟੀ} ਬੰਨ੍ਹਦੇ ਹਨ ਤੇ ਉਨ੍ਹਾਂ
ਬੇਢੰਗਿਆਂ ਨੂੰ ਫੜਦੇ ਹਨ ਜੋ ਫੜੇ ਜਾਣ ਦਿੰਦੇ ਹਨ।

%%K t02-28-1 p
28. --- ਮੈਦਾਨ ਦੇ ਵਿਚਕਾਰ, ਇਹ ਰਹੀ ਇਕ ਬਹੁਤ ਫ਼ਰਾਂਸੀਸੀ ਖੇਡ, ਭਾਵੇਂ ਕੁਝ ਖ਼ੁਰਦਰੀ:
ਇਹ \VUgras{ਰਿੱਛ} \textsuperscript{(32)} ਹੈ, ਜੋ \VUgras{ਪਤੰਗ} \textsuperscript{(33)}
ਦੀ ਸ਼ਾਂਤ ਖੇਡ ਤੋਂ, ਇੱਥੋਂ ਤਕ ਕਿ ਅੰਗਰੇਜ਼ੀ ਖੇਡ \VUgras{ਟੈਨਿਸ} \textsuperscript{(34)}
ਤੋਂ ਵੀ, ਕਿਤੇ ਵੱਧ ਰੌਲੇ ਵਾਲੀ ਹੈ।

%%K t02-29-1 p
29. --- ਉਹ ਆਖ਼ਰੀ ਖੇਡ ਵੱਡੇ ਵਿਦਿਆਰਥੀਆਂ ਦਾ ਆਨੰਦ ਹੈ। ਸਾਡੇ ਅਧਿਆਪਕ ਆਪ ਵੀ ਉਸ ਨੂੰ
ਖ਼ੁਸ਼ੀ ਨਾਲ ਅਪਣਾਉਂਦੇ ਹਨ। ਉਸ ਵਿਚ ਬੜੀ ਸਿਆਣਪ ਤੇ ਫੁਰਤੀ ਚਾਹੀਦੀ ਹੈ, ਤੇ ਉਹ ਮੈਨੂੰ ਬਹੁਤ
ਚੰਗੀ ਲਗਦੀ ਹੈ। \VUgras{ਪੀਂਘ} \textsuperscript{(35)} ਮੈਂ ਕੁੜੀਆਂ ਲਈ ਛੱਡ ਦਿੰਦਾ ਹਾਂ।

%%K t02-30-1 p
30. --- ਇਕ ਹੋਰ ਅੰਗਰੇਜ਼ੀ ਖੇਡ ਵੀ ਮੈਨੂੰ ਚੰਗੀ ਨਹੀਂ ਲਗਦੀ, ਜੋ ਸੌਖੀ ਹੀ ਖ਼ਤਰਨਾਕ ਹੋ
ਸਕਦੀ ਹੈ।

%%K t02-30-1b p
ਮੇਰਾ ਮਤਲਬ \VUgras{ਫੁੱਟਬਾਲ} \textsuperscript{(36)} ਤੋਂ ਹੈ, ਜੋ ਮੇਰੇ ਵੱਡੇ ਭਰਾ ਦਾ
ਆਨੰਦ ਹੈ। ਮੈਨੂੰ ਸਾਡੀ ਪੁਰਾਣੀ ਖੇਡ \VUgras{ਕੈਦੀ-ਅੱਡਾ} \textsuperscript{(38)} ਵੱਧ
ਚੰਗੀ ਲਗਦੀ ਹੈ, ਓਨੀ ਹੀ ਸਿਹਤ ਵਾਲੀ ਤੇ ਕਿਤੇ ਘੱਟ ਸਖ਼ਤ।

%%K t02-tit-7 sub
\VUsaut{4.6mm}
\VUpk{10.2pt}{40}{ਸਾਈਕਲ ਤੇ ਗੇਂਦ ਦੀਆਂ ਖੇਡਾਂ।}
\VUsaut{3.0mm}

%%K t02-31-1 p
31. --- ਮੈਨੂੰ ਮੈਦਾਨ ਦੇ ਚੁਫੇਰੇ \VUgras{ਸਾਈਕਲ} \textsuperscript{(37)} ਚਲਾਉਣਾ ਵੀ
ਚੰਗਾ ਲਗਦਾ ਹੈ।

%%K t02-32-1 p
32. --- ਅਗਲੇ ਵਰ੍ਹੇ ਮੈਂ \VUgras{ਘੋੜਸਵਾਰੀ} \textsuperscript{(39)} ਸਿੱਖਾਂਗਾ। ਘੋੜਾ
ਕਿੰਨਾ ਸੋਹਣਾ ਲਗਦਾ ਹੈ ਜਦੋਂ ਉਹ ਠੁਮਕਦਾ ਜਾਂ ਸੋਹਣੀ ਦੁਲਕੀ ਤੁਰਦਾ ਹੈ, ਤੇ ਸਰਪਟ ਰੁਕਾਵਟਾਂ
ਟੱਪ ਜਾਂਦਾ ਹੈ!

%%K t02-33-1 p
33. --- ਗਰਮੀਆਂ ਵਿਚ ਅਸੀਂ \VUgras{ਤੈਰਨਾ} \textsuperscript{(40)} ਸਿੱਖਦੇ ਹਾਂ। ਅਸੀਂ
ਦਰਿਆ ਦੇ ਸਾਫ਼ ਠੰਢੇ ਪਾਣੀ ਵਿਚ ਆਨੰਦ ਨਾਲ ਚੁੱਭੀ ਲਾਉਂਦੇ ਤੇ ਨ੍ਹਾਉਂਦੇ ਹਾਂ। ਕਦੇ ਕਦੇ ਅੰਤ
ਵਿਚ ਅਸੀਂ ਇਕ ਕਾਗਜ਼ ਦਾ \VUgras{ਗੁਬਾਰਾ} \textsuperscript{(41)} ਉਡਾਉਂਦੇ ਹਾਂ, ਜੋ ਗਰਮ
ਹਵਾ ਨਾਲ ਭਰਦੇ ਹੀ ਗੰਭੀਰਤਾ ਨਾਲ ਉੱਤੇ ਚੜ੍ਹ ਜਾਂਦਾ ਹੈ।

%%K t02-34-1 p
34. --- ਹੋਰ ਵੀ ਬਹੁਤ ਸਾਰੀਆਂ ਖੇਡਾਂ ਹਨ, ਪਰ ਉਨ੍ਹਾਂ ਨੂੰ ਗਿਣਦੇ ਗਿਣਦੇ ਮੈਂ ਕਦੇ ਨਾ
ਮੁੱਕਾਂਗਾ, ਤੇ ਇਸ ਨਿੱਕੀ ਜਿਹੀ ਕਹਾਣੀ ਤੋਂ ਹੀ ਵਿਖਾਈ ਦਿੰਦਾ ਹੈ ਕਿ ਸਾਡੇ ਸੋਹਣੇ ਮੈਦਾਨ
ਉੱਤੇ ਵੀਰਵਾਰ ਨੀਰਸ ਹੋਣ ਤੋਂ ਬਹੁਤ ਦੂਰ ਹਨ।

%%K t02-34-2 p
ਨੋਟ। --- \textit{ਅਧਿਆਪਕ ਇਸ ਅਧਿਆਇ ਨੂੰ ਸੌਖਾ ਹੀ ਪੂਰਾ ਕਰ ਸਕਣਗੇ, ਹਰ ਅਹਿਮ ਖੇਡ ਦਾ
ਵੱਧ ਵਿਸਤਾਰ ਨਾਲ ਵੇਰਵਾ ਦੇ ਕੇ।}
\VUsaut{6.0mm}
\VUfilet{22mm}
